\subsection{Vanishing finite common-character tests with escaping support}

\begin{statement}{A.133}{Proved}
{A completed harmonic tuple with no fixed-centre new coordinate}
There are three classes \(A_0,A_1,A_h\in\widehat{\Q^\times}\),
where \(h=1/2\), and ordinary centres
\[
 t_p\in\Q_p\setminus\{0,1,h,2,-1\}
                         \quad(p<\infty)
\tag{A.133.1}
\]
such that their actual localizations are
\[
 (A_0)_p=\kappa_p(t_p),\qquad
 (A_1)_p=\kappa_p(t_p-1),\qquad
 (A_h)_p=\kappa_p(t_p-h).
\tag{A.133.2}
\]
Their reductions modulo four and their real signs are
those of
\[
       R=\frac14,\qquad
       (a_0,a_1,a_h)=\left(\frac14,-\frac34,-\frac14\right);
                  \qquad t_\infty=R.
\tag{A.133.3}
\]
Each \(A_i\) has infinitely many nonzero ordinary integer
prime valuations.

For the new puncture \(c=2\), every finite-exponent
common-character obstruction of A.128 is trivial.
More precisely, at every \(4\mid N\), choose rational
representatives of \(A_i\bmod N\) and form the old cover
\(C_N:T-i=a_i(N)X_i^N\).
Its actual boundary common field is always
\[
                  B_{N,2}=K_2=\Q(\sqrt2,\sqrt{-3}),
                    \qquad \gamma_{C_N,2}=1.
\tag{A.133.4}
\]
For every finite set of places, the prescribed old
points with centres \(t_p\), and with centre \(R\) at
the real place when prescribed, lift to an everywhere
locally soluble augmented cover
\(T-2=bY^N\) for some \(b\in\Q^\times\).
This finite controlled-point assertion holds at every
positive exponent \(N\), by the explicit power maps.

Nevertheless there is no rational squareclass whose
localizations are \([t_p-2]\) at all finite primes,
and hence no completed global coordinate for this
one fixed family \(t_p-2\).
The construction gives infinitely many distinct primes
\(p_j\) with the literal equalities
\[
                     t_{p_j}=2+p_j,\qquad
                     v_{p_j}(t_{p_j}-2)=1.
\tag{A.133.5}
\]
Thus every finite prescribed-point test succeeds,
while the all-place tuple fails the restricted-product
condition of A.132 already at exponent two.
\end{statement}
\begin{proof}
We give the exact refinement and local-point choices.
They use the mixed-exponent root-stabilizer argument
proved in (A.90.3)--(A.90.10); the calculations below
verify its hypotheses for these new coefficients.

\emph{The four fixed quadratic fields.}
For each old boundary \(i\in\{0,1,h\}\), put
\[
 B_i=\Q\left(\sqrt{(i-k)/a_k}:k\in\{0,1,h\},\ k\ne i\right).
\]
The actual radicands are
\[
\begin{array}{c|cc|c}
 i&\multicolumn{2}{c|}{(i-k)/a_k,\ k\ne i}&B_i\\ \hline
 0&4/3&2&\Q(\sqrt2,\sqrt3)\\
 1&4&-2&\Q(\sqrt{-2})\\
 h&2&2/3&\Q(\sqrt2,\sqrt3).
\end{array}
\tag{A.133.6}
\]
Neither real field contains \(i=\sqrt{-1}\), and the
squareclass span \(\{1,-2\}\) in the middle row does
not contain \(-1\).
At the new boundary \(2\), the three radicands are
\[
           (D_0,D_1,D_h)=(8,-4/3,-6).
\tag{A.133.7}
\]
Their squareclass span is exactly
\(\{1,2,-3,-6\}\), again excluding \(-1\).
This proves the displayed value of \(K_2\) in
(A.133.4).

All refinements below multiply each coefficient by
a fourth power in \(\Q^\times\).
Consequently these four quadratic fields remain
exactly the displayed fields.
For arbitrary mixed exponents \(4\mid n_i\), the
actual boundary algebras have equations
\(Z_k^{n_k}=(i-k)/a_k'\), \(k\ne i\), or
\(Z_k^{n_k}=(2-k)/a_k'\), respectively.
The proof (A.90.3)--(A.90.10), including its odd-part
CRT correction, says their common abelian fields
are respectively \(B_i\) and \(K_2\).
This applies to the mixed exponents themselves,
without replacing their fibre maps by an equal-exponent
cover.

\emph{The controlled covers and their invariants.}
Start with
\[
 C_0:\quad T-i=a_iX_i^4\quad(i=0,1,h),\qquad
                      X_0X_1X_h\ne0,\quad T\ne2,-1,
\tag{A.133.8}
\]
and its rational point \((R,1,1,1)\).
At a stage use the actual open curve
\[
 C_j:\quad T-i=a_{i,j}X_i^{n_{i,j}},
       \qquad 4\mid n_{i,j},\quad
                X_0X_1X_h\ne0,\quad T\ne2,-1.
\tag{A.133.9}
\]
A finite controlled set \(S_j\) contains all
coefficient, exponent, puncture-difference and
denominator primes, and the ramified primes of
the fixed \(B_i\).
At each \(p\in S_j\) retain an actual point of \(C_j\).
Its coordinate \(T\) will remain fixed thereafter.
All retained root coordinates are units and satisfy
\[
                  (X_i,\chi)=0
       \quad\text{for every character \(\chi\) of \(B_i\)}.
\tag{A.133.10}
\]
At the real place retain \(T=R\) and positive root
coordinates.
Outside \(S_j\), require local points with all root
coordinates units and both \(T-2\) and \(T+1\) units.
These assertions hold initially after enlarging
\(S_0\), using the rational point: its roots are one,
and \(R-2=-7/4\), \(R+1=5/4\) are units away
from two, five and seven.
Include those primes in \(S_0\), together with all
prime supports of the nonzero differences \(i+1\)
for \(i=0,1,h\).

We will use a uniform good-prime bound at each step.
For the three exponents put \(D=\prod_i n_i\) and
\(N=\operatorname{lcm}_i n_i\).
The geometric degree is \(D\), the finite inertia
orders are \(n_i\), and the order at infinity is \(N\).
The same root-label maps as in (A.90.18) give
\[
 2g-2=-2D+\sum_iD(1-1/n_i)+D(1-1/N),\qquad
 b=\sum_iD/n_i+D/N+2D.
\tag{A.133.11}
\]
The last \(2D\) points remove the two unramified fibres
\(T=2,-1\).
Thus \(2g\le2D+2\), \(b\le6D\), and
\(B(\boldsymbol n)=(8D+8)^2\) is a sufficient
strict bound for \(p+1-2g\sqrt p-b>0\).
At a larger prime with all coefficients and
puncture differences units and \(p\nmid N\),
the good-reduction point count and Hensel lifting
supply the required unit-root point with both \(T-2\)
and \(T+1\) units, by the exact argument of A.90.

\emph{Adding one forced centre.}
Before each refinement choose a new prime \(p_j\)
outside every controlled set and every prime used
earlier.  Let \(F_j\) be the compositum of the
splitting field of
\[
          Z_i^{n_{i,j}}=(2-i)/a_{i,j}\quad(i=0,1,h)
\tag{A.133.12}
\]
with the three fixed fields \(B_i\).
Chebotarev chooses \(p_j>3\) splitting completely in
\(F_j\), larger than the current bounds and not
dividing any exponent, coefficient or difference.
Its residue field contains a nonzero root of
each equation (A.133.12).
For the prescribed rational centre \(T=2+p_j\),
the polynomials
\[
             a_{i,j}X_i^{n_{i,j}}-(2+p_j-i)
\tag{A.133.13}
\]
have those same roots modulo \(p_j\), with unit
derivatives \(n_{i,j}a_{i,j}X_i^{n_{i,j}-1}\).
Hensel lifting therefore gives actual unit
coordinates in \(\Q_{p_j}\).
Every character of every \(B_i\) is locally
trivial at this prime, so (A.133.10) holds.
Add this point to the controlled set and retain
its \(T\) permanently.
The equality \(v_{p_j}(T-2)=1\) is already exact.
Also \(T+1=3+p_j\) is a unit at \(p_j\).
The prime can be chosen distinct at every step,
so the procedure gives infinitely many such primes.

\emph{The exact old-coordinate refinement.}
Choose one old index \(i\) and an integer \(h'\ge2\).
We refine it by the actual morphism
\[
\begin{aligned}
 C_{j+1}&\longrightarrow C_j,\\
 X_i&=uY_i^{h'},\qquad X_k=Y_k\ (k\ne i),\qquad T=T,\\
 n_i'&=h'n_i,\qquad a_i'=a_i u^{n_i},
       \qquad n_k'=n_k,\ a_k'=a_k\ (k\ne i).
\end{aligned}
\tag{A.133.14}
\]
Here \(u>0\) is the rational number now constructed;
the unsubscripted coefficients in this paragraph
are the current ones.
The equations verify the map by substitution.

Enlarge the controlled set to \(S\) containing the
new bound (A.133.11), the new exponent primes,
all primes at most \(j+1\), and the ramification
of the current boundary splitting field at \(T=i\).
At newly controlled primes choose the unit-root
points with \(T-2\) and \(T+1\) units supplied by the invariant.
Their symbols (A.133.10) vanish because the \(B_k\)
are unramified and their root coordinates are units.
All previous points, including the forced centres,
are retained.

Let \(x_p\) be the current \(X_i\)-coordinate at
each \(p\in S\).  It is a unit.
Choose precisions \(M_p>2v_p(h')+4\), enlarged so
that \(\mathcal R=\Q(\mu_M)\) contains \(B_i\), where
\(M=\prod_{p\in S}p^{M_p}\).
Let \(\rho\in(\Z/M)^\times\) have
\(\rho\equiv x_p\pmod{p^{M_p}}\) for all \(p\in S\).
The individual zero symbols (A.133.10), the positive
real root, and local cyclotomic reciprocity give
\[
          \rho^{-1}|_{B_i}
             =\prod_{p\in S}\operatorname{rec}_p(x_p)|_{B_i}
             =1.
\tag{A.133.15}
\]
There is no nonunit valuation factor in this instance;
all old root coordinates have been kept units.

The actual residue algebra at \(X_i=0\) is
\[
 \mathcal E_i=
   \Q[Z_k^{\pm1}:k\ne i]/
          (Z_k^{n_k}-(i-k)/a_k:k\ne i).
\tag{A.133.16}
\]
Its common abelian field is \(B_i\), as proved above.
The images in \(\operatorname{Gal}(\mathcal R/\Q)\)
of all its root-tuple stabilizers therefore generate
\(\operatorname{Gal}(\mathcal R/B_i)\), by the exact
restriction and Galois-correspondence equality
(A.90.25).
Write \(\rho\) as a product of these images.
Chebotarev supplies distinct fresh primes \(q_l\)
with those restrictions and actual fixed residue
root tuples, avoiding \(S\) and all previous primes.
Choose one more prime \(q_*\) splitting the whole
compositum of the boundary splitting field and
\(\mathcal R\), and put
\[
                         u=q_*\prod_lq_l>0.
\tag{A.133.17}
\]
Then \(u/x_p\in1+p^{M_p}\Z_p\) for all \(p\in S\),
and \(v_{q_*}(u)=1\).

At \(p\in S\), retain \(T\) and every other root,
and choose
\[
                    Y_i=(x_p/u)^{1/h'}\in1+p^3\Z_p.
\tag{A.133.18}
\]
The stated precision gives this root by Hensel's
lemma.  It is a square, including when \(p=2\);
therefore its symbols against the quadratic
characters of \(B_i\) vanish.
The other roots and fixed \(B_k\) are unchanged.
At the real place take the positive root, which
retains \(T=R\).
Every forced centre \(2+p_j\) is unchanged by this
literal morphism.

At a fresh coefficient prime \(q\mid u\), take
\(Y_i=1\) and \(T=i+a_i u^{n_i}\).
The other root equations reduce to the selected
nonzero boundary roots and have unit derivatives.
They lift to unit roots.
Every \(B_k\) is unramified there, so all root
symbols vanish.
Moreover both \(T-2\) and \(T+1\) are units:
\(i-2\) and \(i+1\) are units at this fresh prime
and \(a_i u^{n_i}\) has positive valuation.
At every other prime outside \(S\) the new good-prime
bound supplies the unit-root point with \(T-2\)
and \(T+1\) units.
Adding the fresh coefficient primes to the controlled
set proves all the invariants for the next step.

\emph{The completed tuple and the infinite supports.}
Run cycles \(h'=2,3,4,\ldots\), refining the three
old coordinates in order during each cycle.
After cycle \(h'\), their common exponent is
\(4h'!\), cofinal under divisibility by every positive
integer.  Equation (A.133.14) changes a coefficient
by a power of its old exponent.
Thus the rational coefficient classes define
compatible global elements \(A_i\in\widehat{\Q^\times}\),
with reductions (A.133.3) modulo four.
Every finite prime eventually becomes controlled.
Thereafter \(T=t_p\) is fixed, and at all sufficiently
large exponents the actual local equation gives
\[
                [a_{i,j}]=[t_p-i]
                        \quad\text{in }\Q_p^\times/(\Q_p^\times)^N.
\tag{A.133.19}
\]
This proves (A.133.2) in the full local completion.
It also proves that its values are ordinary nonzero
numbers, not just a compatible collection of varying
local centres.  Positivity of every \(u\) gives the
stated global signs and real point.

At each refinement of coordinate \(i\), the fresh
prime \(q_*\) gives valuation \(n_i>0\) in its new
coefficient and valuation zero in the other two.
Its retained new roots are units.
Every subsequent refinement therefore has unit
factor \(u\) there, so those integer valuations
remain unchanged.
Infinitely many refinements of every coordinate
prove the infinite-support assertion.
The separately chosen primes \(p_j\) remain
distinct and retain exactly (A.133.5).

\emph{Every finite controlled initialization succeeds.}
At \(4\mid N\), take representatives from a
sufficiently late completed stage.
Their ratios to the initial \(a_i\) are fourth
powers, so the three boundary squareclasses at
two are still (A.133.7).
The root-stabilizer result A.80 therefore gives
the exact common field \(B_{N,2}=K_2\).
The old finite cover has the actual ordinary local
points supplied by (A.133.2), and a real point by
the retained signs.
Its individual character identities are
\[
             (T-2,D_i)=(a_i(N),D_i)
                                  \quad\text{in }\operatorname{Br}(C_N),
\tag{A.133.20}
\]
by the literal quadratic norm in (A.128.27).
These characters generate all characters of \(K_2\).
Global reciprocity for the constants on the right
proves \(\gamma_{C_N,2}=1\).
Alternatively the same vanishing is the exact
hypothesis checked in A.83 by its principal-idele
calculation.
That theorem supplies an augmented all-place cover
lifting any prescribed finite list of the actual
old points.
For an arbitrary positive \(N\), apply this result
at \(4N\) and use the actual coordinate-power maps
to exponent \(N\); the local centre is unchanged.
This proves every finite controlled-point assertion.

\emph{Failure for the one fixed all-place tuple.}
If \(b\in\Q^\times\) represented its proposed new
global squareclass, choose one of the infinitely
many \(p_j\) outside the finite prime support of \(b\).
Then \(v_{p_j}(b)=0\), while (A.133.5) gives
\(v_{p_j}(t_{p_j}-2)=1\).
Their quotient cannot be a square in \(\Q_{p_j}\).
This excludes every rational squareclass and hence
every completed new coordinate.
It is exactly the failure of (A.132.5) at \(N=2\).
The common-character obstructions in (A.133.4)
remain zero because they test the finite-cover
initialization with freely chosen outside points,
through the restriction maps (A.132.9).
The assertion concerns completed Kummer coordinates;
no fundamental-group section is inferred from it.

\emph{An effective convention specifying one tuple.}
Fix the increasing order of rational primes and
lexicographic orders on every finite residue set
and on the finite root-label groups in
(A.90.4)--(A.90.9).
At each forced-centre step take the least admissible
split prime after the specified exclusions, then
the first nonzero residue root in each coordinate
and its unique simple-root Hensel branch.
At a newly controlled good prime take the first
interior residue point of the open curve counted
in (A.133.11), fix the least nonnegative integral
lift of its \(T\)-coordinate, and take the unique
Hensel branches of its nonzero roots.
At every coefficient refinement compute the
actual mixed-exponent stabilizers by
(A.90.5)--(A.90.9), enumerate words in their
finite restriction images, and take the first
word whose product is \(\rho\).
For its successive Frobenius classes take the
least admissible distinct primes, and take the
least further completely split prime for \(q_*\).
The close-to-one root in (A.133.18) is its specified
Hensel branch.
Use the least integer precisions meeting all the
stated inequalities and conductor requirements.
The point-count, stabilizer-generation and
Chebotarev proofs show that each finite selection
terminates.  These conventions, together with
the stated factorial cycle, specify one actual
infinite tuple.  They do not assert that its
infinite prime recursion has been executed.
\end{proof}

The retained exact certificate
\path{code/fixed_local_escaping_support_certificate.py}
checks the four fixed squareclass fields, 576 literal
coefficient and forgetful-map identities, 512 mixed
genus bounds, factorial cofinality, and fifteen actual
split-prime fixed-centre examples.  It executes 180
Hensel lifts through modulus \(p^5\), together with
the old-root precision and valuation exclusions.
Its saved output states the primes used and makes
no claim to have executed the infinite Chebotarev
recursion.

